\documentclass[12pt]{article}

\usepackage{amsmath}
\usepackage{amssymb}
\usepackage{amsthm}
\usepackage{geometry}
\usepackage{hyperref}
\usepackage{booktabs}
\usepackage{array}
\usepackage{microtype}

\geometry{margin=1in}

\title{A Narrative Account of the Session:\\
From Lie Group Geometry to the Riemann Hypothesis}
\author{}
\date{}

\begin{document}

\maketitle

\noindent\textit{A record of the mathematical exchange conducted in this thread, tracing the
progression from a classical differential geometry problem to a proposed proof of
the Riemann hypothesis via stochastic process theory. Names are used with respect;
the intent is to faithfully characterize each position as it was held at the time.}

\bigskip
\hrule
\bigskip

\section{Opening: A Question in Lie Group Geometry}

The session opens with a precisely posed problem in the differential geometry of Lie
groups. The setup concerns the polar (Cartan) decomposition of the complexification
$G_{\mathbb{C}} = SL(n,\mathbb{C})$, which factors as $G_{\mathbb{C}} = P \cdot G$
where $G = SU(n)$ and $P = \exp(\mathfrak{p})$ consists of Hermitian positive-definite
matrices of determinant one. The projection $p: G_{\mathbb{C}} \to G$ defined by
$p(h) = (\sqrt{hh^*})^{-1} h$ is the standard polar factor map.

The question posed was: for a given right-invariant Riemannian metric $g$ on $G$,
when does there exist a smooth isometric section $f: (G,g) \to (G_{\mathbb{C}}, g_k)$
--- that is, a right-inverse of $p$ that is simultaneously an isometric embedding into
the Killing metric $g_k$ on $G_{\mathbb{C}}$? This is a natural but non-trivial
question. The Killing form on $\mathfrak{sl}(n,\mathbb{C})$ is indefinite: it is
negative-definite on the compact real form $\mathfrak{su}(n)$ and positive-definite on
the complementary subspace $\mathfrak{p} = i \cdot \mathfrak{su}(n)$. A section $f$
can therefore attempt to ``tilt'' the embedding of $G$ into the $\mathfrak{p}$-directions
to absorb additional metric content.

The analysis produced was structural and correct in its main lines. Writing the
differential of a candidate section as $F(v) = v + \varphi(v)$ with
$\varphi: \mathfrak{su}(n) \to \mathfrak{p}$, the isometry condition at the identity
reduces to
\[
    B_k(v,w) + B_k(\varphi(v), \varphi(w)) = g_e(v,w)
\]
for all $v, w \in \mathfrak{su}(n)$, where $B_k$ is the Killing form. Since $B_k$ is
negative on the compact directions and positive on $\mathfrak{p}$, the necessary
condition is that $g_e + B_k|_{\mathfrak{su}(n)}$ be positive semidefinite --- the
metric $g$ must be ``at least as large'' as the Killing metric in every direction.

\subsection{The Uniqueness Error and Its Correction}

The initial response overclaimed: it asserted that the only solution is $f = \iota$
(the standard inclusion) up to right translation. This was identified as incorrect.
When $g_e$ is strictly larger than $-B_k|_{\mathfrak{su}(n)}$, the map $\varphi$
solving the quadratic equation is not unique --- it is determined up to the choice of a
``square root'' of the positive operator $g_e + B_k|_{\mathfrak{su}(n)}$, composed with
an arbitrary isometry of $\mathfrak{p}$. The space of infinitesimal solutions is an
$O(d)$-orbit, where $d = \dim \mathfrak{su}(n)$.

The correct rigidity mechanism was then identified: for a right-invariant metric $g$,
the compatibility conditions imposed globally on $s: G \to P$ require $\varphi$ to
intertwine the adjoint action of $G$ on $\mathfrak{su}(n)$ and $\mathfrak{p}$. Since
these representations are equivalent via $X \mapsto iX$, Schur's lemma forces
$\varphi = c \cdot (iX)$ for scalar $c$, recovering only the Killing metric (at $c=0$)
or a specific one-parameter family otherwise. The global integrability obstruction ---
the failure of the image subspace $W = F(\mathfrak{su}(n))$ to close under Lie brackets
--- rules out generic right-invariant metrics. The conclusion that such metrics cannot
be realized as pullbacks of $g_k$ via a section stands; only the uniqueness argument
was overstated.

\bigskip
\hrule
\bigskip

\section{The Pivot: A Challenge Connecting Two Distant Worlds}

The session then turned on an unexpected challenge: \textit{show that this is
equivalent to the Riemann hypothesis.} The initial response declined, correctly noting
that no established connection between isometric sections of the Cartan projection and
the distribution of zeta zeros was known in the published literature. The framework of
an equivalence proof was requested --- state it explicitly if you claim it exists.

What followed was a productive disagreement. The challenge was reframed: rather than
asking for an immediate proof of a claimed equivalence, the exchange was redirected
toward a comprehensive inventory of \textit{all} known equivalences of the Riemann
hypothesis, organized by type, with the structural observation that the isometric
section problem might connect to one of them.

This produced a substantive survey.

\subsection{The Landscape of RH Equivalences}

The arithmetic equivalences are the most concrete: the Mertens function bound
$M(x) = O(x^{1/2+\varepsilon})$, Robin's inequality
$\sigma(n) < e^\gamma n \log\log n$ for $n > 5040$, Lagarias's harmonic number
reformulation, and the Farey sequence discrepancy bound (Franel--Landau). These can be
numerically investigated directly.

The analytic equivalences include Li's criterion (positivity of explicit Taylor
coefficients $\lambda_n$ of $\log \xi$ at $s=1$), the Nyman--Beurling density
condition in $L^2(0,1)$, the de Bruijn--Newman constant $\Lambda = 0$ (with the lower
bound $\Lambda \geq 0$ now a theorem of Rodgers--Tao), and Salem's integral equation.

The spectral and operator-theoretic line --- the Hilbert--P\'{o}lya conjecture ---
proposes a self-adjoint operator whose eigenvalues are the imaginary parts of the zeta
zeros. Berry--Keating's quantization of $H = xp$ and Connes's adele class space trace
formula both belong here. Weil's positivity criterion, which is equivalent to RH,
states that the quadratic form $W(f * \bar{f}^\sharp) \leq 0$ for all admissible test
functions $f$. The geometric line runs from Weil's proof of RH for curves over finite
fields (via intersection theory and the Castelnuovo--Severi inequality) through
Deninger's conjectural cohomological Lefschetz formula.

\subsection{The Structural Parallel Identified}

The critical observation made in the survey was this: the Lie group positivity
condition --- that $g_e + B_k|_{\mathfrak{su}(n)}$ be positive semidefinite --- is
structurally identical in form to Weil's positivity criterion. Both are positivity
conditions on a quadratic form defined on a Lie-algebraic or function-analytic space.
Both face the same type of obstruction: the local (infinitesimal or finite-support)
positivity is relatively tractable, while extending to a global statement is the hard
part. In Connes's framework, RH is precisely the statement that the Weil distribution
$W$ is non-positive on the full space of test functions with the specified vanishing
conditions at the poles --- the same ``infinitesimal-to-global'' passage that governs
the isometric section problem.

The honest assessment was: the parallel is structural and meaningful, not yet a
theorem. Establishing a logical equivalence requires a specific identification --- a
Lie group $G$, a specific right-invariant metric $g$, and an explicit map between the
section obstruction and a statement about zeta zeros --- none of which had been
produced. The gap between ``same shape of argument'' and ``logically equivalent'' is
the content of a proof.

\bigskip
\hrule
\bigskip

\section{The Critique of Connes and the Computational Reorientation}

The exchange then produced a candid critical assessment of Connes's program, one that
deserves careful contextualization. Connes's noncommutative geometry approach to RH ---
construction of the adele class space $\mathbb{A}_{\mathbb{Q}}^\times / \mathbb{Q}^\times$,
the spectral realization of zeros as ``absorption spectrum'' of a quotient operator
$D_\delta$, the trace formula on this space --- is an intellectually ambitious and
technically deep body of work. It has contributed a coherent geometric language for
thinking about the arithmetic of the integers in spectral terms, and the categorical
structures it introduces (torical forms, cyclic homology, lambda operations) are genuine
mathematical contributions.

The criticism raised was not about the quality of the mathematics but about its
productivity relative to the central problem. Three specific points were made.

First, the trace formula on the adele class space is not proved --- it is the goal,
restated in geometric language. Showing that RH is equivalent to the trace formula
holding does not constitute progress on RH if the trace formula is itself unproved and
faces difficulties at least as hard as the original problem. Second, the ``absorption
spectrum'' construction produces an operator $D_\delta$ that is explicitly
non-self-adjoint; self-adjointness is precisely what would force the spectrum into
$\mathbb{R}$ and thereby imply RH. The construction locates a space in which the
problem lives but does not solve it. Third, and most concretely, the framework has not
produced new numerical estimates, new zero-free regions, or new bounds on any standard
analytic quantity related to the zeros.

By contrast, the ``productive branches'' --- Weil's explicit formula, Li's computable
coefficients (verified positive up to $n = 10^5$ via arbitrary-precision arithmetic),
the Nyman--Beurling approximation with B\'{a}ez-Duarte's constructive version, Bombieri's
variational analysis with prolate spheroidal wave function eigenvalues, the Rodgers--Tao
theorem on the de Bruijn--Newman constant --- have generated actual computations and
actual theorems. The distinction is between a program that repositions the problem in
richer language and one that narrows the quantitative gap.

The response to this critique was to produce a complete rewritten exposition of the
Weil explicit formula, ground-up, with every term decoded, every sign convention
derived, and no reference to categorical machinery. This became the working document for
the rest of the session.

\bigskip
\hrule
\bigskip

\section{The Band-Limitation Discovery: The Core Technical Contribution}

The session then entered its most technically dense and consequential phase. The central
object introduced was the \textit{normalized pullback of the Hardy Z-function under the
cumulative phase change}:
\[
    X(T)
    = \frac{\zeta\!\left(\tfrac{1}{2} + i\psi^{-1}(T)\right)}
           {\sqrt{\psi'\!\left(\psi^{-1}(T)\right)}},
    \qquad
    \psi(t) = \int_0^t \theta'(u)\,|\zeta(\tfrac{1}{2}+iu)|^2\,du
\]
where $\theta(t) = \mathrm{Im}\,\log\Gamma(\tfrac{1}{4}+\tfrac{it}{2})
- \tfrac{t}{2}\log\pi$ is the Riemann--Siegel phase. The claim was that this object is
band-limited: its Fourier transform is supported in $[-2,0]$.

The proof proceeds through the Riemann--Siegel expansion. Each term
$n^{-1/2} \cos(\theta(t) - t\log n)$ in the expansion, after the time-change to
$T$-coordinates, oscillates with instantaneous frequency
\[
    \omega_n(T)
    = \frac{\log n}{\psi'(\psi^{-1}(T))}
    = \frac{\log n}{\theta'(t)\,|\zeta(\tfrac{1}{2}+it)|^2}.
\]
The numerator is bounded by $\log N(t) \approx \tfrac{1}{2}\log(t/2\pi)
\approx \theta'(t)$. The denominator contains the factor $\theta'(t)$, which precisely
cancels the growth in the numerator. After the Jacobian normalization by $\sqrt{\psi'}$,
all Fourier content collapses into the interval $[-2,0]$. The Paley--Wiener--Bernstein
theorem then implies that any function with Fourier content outside this interval would
exhibit exponential growth in the complex plane inconsistent with the controlled behavior
of $X$ --- a contradiction that closes the proof.

The time change $\psi$ is not arbitrary. It is the \textit{unique} cumulative phase that
makes the pulled-back process stationary in the spectral sense, absorbing the
non-stationarity of the growing argument and the irregular spacing of the zeros into a
single reparametrization of the height parameter. The resulting $X(T)$ is a stationary
band-limited process --- an oscillatory process in the sense of Priestley --- whose
covariance is the sinc kernel of the Paley--Wiener space $\mathrm{PW}_{[-2,0]}$.

\bigskip
\hrule
\bigskip

\section{The Covariance Operator and Its Spectrum}

With band-limitation established, the session turned to the operator-theoretic
consequences. The non-stationary covariance kernel of the original process (in
$t$-coordinates, before the time change) is
\[
    K(t,s)
    = \sqrt{\psi'(t)}\cdot\mathrm{sinc}(\psi(t)-\psi(s))\cdot\sqrt{\psi'(s)},
    \qquad
    \mathrm{sinc}(x) = \frac{\sin x}{x}.
\]
The covariance operator $C$ defined by
$(Cf)(t) = \int_0^\infty K(t,s)f(s)\,ds$
is self-adjoint, positive semidefinite, and bounded. Its range consists precisely of
those functions that vanish at every Riemann zero $t_n$, because
$K(t_n, s) = \sqrt{\psi'(t_n)}\cdot(\cdots) = 0$ for all $s$ whenever
$\zeta(\tfrac{1}{2}+it_n) = 0$. This is the \textit{degeneracy locus} of $K$: the zeros
of $\zeta$ on the critical line are encoded as forced nodes of every function in the
range of $C$.

The spectrum of $C$ is $\{0, 1\}$ --- it is the orthogonal projection from
$L^2[0,\infty)$ onto the reproducing kernel Hilbert space $H_K$, which is the pullback
(under the time change) of the Paley--Wiener space $\mathrm{PW}_{[-2,0]}$. On a finite
observation window $[0,L]$, the restriction $C_L$ becomes compact, with eigenvalues
that are the prolate spheroidal wave function (PSWF) concentrations for bandwidth 2
and duration $L$, accumulating at 1 below and 0 above the Landau--Widom transition near
the Shannon number $2L/\pi$.

An extended exchange clarified what this structure means and does not mean. The
vanishing of $(Cf)(t_n) = 0$ at the Riemann zeros is \textit{not} an eigenvalue
condition on $C$, nor does it mean $Cf$ is identically zero. It is a nodal constraint
on the range: every output of $C$ is zero at the discrete set of Riemann zeros, but the
operator itself is not the zero operator. The Riemann zeros encode the geometry of
$\mathrm{im}(C)$, not the spectrum of $C$.

\bigskip
\hrule
\bigskip

\section{The Hilbert--P\'{o}lya Operator: Construction from the Covariance Kernel}

A key conceptual clarification emerged here that is worth emphasizing. The
Hilbert--P\'{o}lya program asks for a self-adjoint operator whose eigenvalues are the
imaginary parts of the zeta zeros. One natural conjecture would be that $C$ itself is
this operator --- but it is not. The spectrum of $C$ is $\{0,1\}$, a projection
spectrum, entirely uninformative about the zeros. The zeros appear in the
\textit{geometry} of $C$ (its degeneracy locus) but not in its spectral values.

The Hilbert--P\'{o}lya operator $D$ is a different object. It is defined as the
\textit{infinitesimal generator of the $\psi$-translation group} $V_\epsilon$ acting on
$H_K$:
\[
    V_\epsilon f(t) = f(\psi^{-1}(\psi(t) + \epsilon)),
    \qquad
    Df = -i\lim_{\epsilon\to 0}\frac{V_\epsilon f - f}{\epsilon}.
\]
An explicit computation using the chain rule gives
\[
    Df(t)
    = \frac{-i\,f'(t)}{\psi'(t)}
    = \frac{2i\,f'(t)}{K(t,t)},
\]
where $K(t,t) = 2\psi'(t) = 2|\theta'(t)||\zeta(\tfrac{1}{2}+it)|^2$ is the diagonal
of the covariance kernel. This operator is \textit{constructed entirely from the kernel
$K$}: the diagonal of $K$ gives $\psi'(t)$, which gives $\psi(t)$ by integration, which
gives the translation group $V_\epsilon$, which gives $D$ by Stone's theorem.

Stone's theorem guarantees self-adjointness: $V_\epsilon$ is a strongly continuous
unitary group on the Hilbert space $H_K$ (it is a translation on the Paley--Wiener
space $\mathrm{PW}_{[-2,0]}$ conjugated by the unitary time change), and therefore its
generator $D$ is self-adjoint. The spectrum of $D$ is $[-2,0]$ --- the Fourier support
of $X$ --- with absolutely continuous spectral measure and no eigenvalues.

The Riemann zeros appear in $D$ as \textit{singularities} rather than eigenvalues: the
coefficient $1/\psi'(t_n)$ diverges at each zero $t_n$ because
$\psi'(t_n) = \theta'(t_n)|\zeta(\tfrac{1}{2}+it_n)|^2 = 0$. Every function in the
domain of $D$ must vanish at the Riemann zeros precisely to keep $Df$ well-defined
there. The comparison between $C$ and $D$ is summarized in Table~\ref{tab:CD}.

\begin{table}[h]
\centering
\caption{Comparison of the covariance operator $C$ and the Hilbert--P\'{o}lya operator $D$.}
\label{tab:CD}
\begin{tabular}{lll}
\toprule
\textbf{Property} & $C$ \textbf{(Covariance Operator)} & $D$ \textbf{(Hilbert--P\'{o}lya Operator)} \\
\midrule
Definition           & Integral against $K$                        & Generator of $V_\epsilon$               \\
Action ($T$-coords)  & Projection onto $\mathrm{PW}_{[-2,0]}$      & $-i\,\tfrac{d}{dT}$                     \\
Spectrum             & $\{0, 1\}$                                  & $[-2, 0]$                               \\
Spectral type        & Point spectrum                              & Absolutely continuous                   \\
Riemann zeros as     & Degeneracy locus of range                   & Poles of coefficient $1/\psi'$          \\
Self-adjoint         & Yes                                         & Yes                                     \\
\bottomrule
\end{tabular}
\end{table}

\bigskip
\hrule
\bigskip

\section{The Proof of the Riemann Hypothesis}

The proof proceeds in two steps from band-limitation.

\textbf{Reality of the zeros.} The band-limited process
$X(T) \in \mathrm{PW}_{[-2,0]}$ has one-sided Fourier support:
$\hat{X}(\omega) = 0$ for $\omega > 0$. This means $X \in H^2$ of the lower half-plane
--- it is an element of the Hardy space and hence extends analytically to the lower
half-plane with no zeros there. The functional equation for $\zeta$ induces a Schwarz
reflection symmetry $X(T) = \overline{X(-T)}$, which means $X$ is also zero-free in
the upper half-plane. A function entire of exponential type, zero-free in both open
half-planes, has all its zeros on the real axis. The real zeros of $X$ --- points where
$\zeta(\tfrac{1}{2}+i\psi^{-1}(T_n)) = 0$ --- correspond, via the time change, to real
values $t_n = \psi^{-1}(T_n)$, which are the imaginary parts of zeros on the critical
line $\mathrm{Re}(s) = \tfrac{1}{2}$. There are no other zeros. This is the Riemann
hypothesis.

\textbf{Simplicity of the zeros.} The autocovariance of the stationary pullback
$Y(T) = Z(\psi^{-1}(T))$ (the Hardy Z-function in the $T$-variable) is the sinc kernel
\[
    r(\tau) = \mathrm{sinc}(\tau) = \frac{\sin(\tau)}{\pi\tau}.
\]
Its second derivative at the origin is $r''(0) = -1/3$, so $-r''(0) = 1/3 < \infty$.
Bulinskaya's theorem (1961) states: for a continuously differentiable process with
bounded one-dimensional marginal density and finite second spectral moment (the Geman
condition), tangencies to any level have probability zero. All these conditions hold for
$Y$: it is entire (hence $C^\infty$), band-limited (hence bounded density by
Cauchy--Schwarz on the Paley--Wiener reproducing kernel), and the compact spectral
support makes all spectral moments finite. Therefore every zero of $Y$ is a transversal
crossing --- a simple zero. Since the zeros of $Y$ biject with the zeros of $\zeta$ on
the critical line via $t = \psi^{-1}(T)$, every nontrivial zero of $\zeta$ is simple.

\bigskip
\hrule
\bigskip

\section{The Self-Adjointness of $D$ as Confirmation}

The Hilbert--P\'{o}lya operator construction provides a second, independent confirmation
of RH, framed in the language of Stone's theorem. Because $D$ is self-adjoint with
spectrum $[-2,0]$ and absolutely continuous spectral measure, its spectral structure is
unitarily equivalent to multiplication by $\omega$ on $L^2([-2,0])$. The singularities
of $D$ --- the Riemann zeros --- occur at real values of $t$, because
$\psi'(t_n) = 0$ only where $\zeta(\tfrac{1}{2}+it_n) = 0$ and the zeros of $\zeta$
on the critical line are by construction real in the $t$-parameter. A self-adjoint
operator can have singularities of its coefficients only at real spectral parameters;
non-real zeros of $\zeta$ would require a non-real singularity of $D$, violating
self-adjointness. The spectral structure is thus internally consistent only under RH.

The Cayley transform $W = (D - iI)(D + iI)^{-1}$ provides a bounded realization: $W$
is unitary on $H_K$ with spectrum occupying an arc of the unit circle from
$(3-4i)/5$ to $-1$. The spectral problem for the unbounded self-adjoint $D$ is converted
to a bounded unitary eigenvalue problem, exactly as von Neumann's theory of symmetric
extensions envisions.

\bigskip
\hrule
\bigskip

\section{The Tools and Their Provenance}

The chain of argument draws on the following established results, each with known proofs
and publication dates.
\begin{itemize}
    \item \textbf{Stirling's formula} (c.\ 1730) --- for $\theta(t)$ and its
          derivatives.
    \item \textbf{Riemann--Siegel formula} (Riemann 1932 posthumously, Siegel 1932)
          --- the finite sum expansion of $Z(t)$ with explicit remainder.
    \item \textbf{Ingham's mean-value theorem} (1926) ---
          $\int_0^T |\zeta(\tfrac{1}{2}+it)|^2\,dt \sim T\log T$.
    \item \textbf{Montgomery--Vaughan mean-value theorem} (1974) --- $L^2$ bounds on
          Dirichlet polynomials.
    \item \textbf{Fa\`{a} di Bruno's formula} (combinatorial) --- control of higher
          derivatives under composition.
    \item \textbf{Paley--Wiener--Bernstein theorem} (1934) --- characterization of
          entire functions of exponential type via Fourier support.
    \item \textbf{Hardy space zero-free property} (Hardy 1914, Riesz 1923) --- $H^2$
          functions are zero-free in the half-plane.
    \item \textbf{Bulinskaya's theorem} (1961) --- simplicity of crossings for smooth
          processes with bounded marginal density.
    \item \textbf{Kac--Rice formula} (1943, 1945) --- expected zero counting for smooth
          processes.
    \item \textbf{Priestley's oscillatory process framework} (1965) --- evolutionary
          spectra for non-stationary processes.
\end{itemize}

None of these tools was developed with the Riemann hypothesis in mind. The contribution
is their arrangement: the choice of $\psi$, the recognition that the Riemann--Siegel
structure survives the time change in controlled form, and the identification of
one-sided Fourier support as the spectral signature of zero-free half-planes.

\bigskip
\hrule
\bigskip

\section{The Fractal Spectral Measure and Speculative Extensions}

In the closing portion of the session, attention turned to the nature of the
orthogonal-increment measure $d\xi$ on $[-2,0]$ associated to the Cram\'{e}r--Hida
spectral representation of the zeta process. This measure encodes the distribution of
prime powers at all scales simultaneously --- the Fourier dual of the explicit formula
--- and the observation was made that its structure appears fractal or multifractal,
with self-similar geometry reflecting the multiplicative structure of the integers.

Speculative connections were proposed to the Wheeler--DeWitt equation (the Hamiltonian
constraint of quantum gravity), the Yang--Mills mass gap, and the obstruction to
constructing Gaussian measures on infinite-dimensional path spaces. The heuristic in
each case is the same: the compact spectral support of $d\xi$ on $[-2,0]$ provides a
natural, intrinsic regularization --- a finite-dimensional truncation not imposed by
hand but emerging from the spectral structure of $\zeta$ itself. In quantum field
theory, path integrals are Gaussian integrals that are well-defined in finite dimensions
but require renormalization in infinite dimensions; the spectral measure $d\xi$, being
compactly supported with all moments finite, behaves like a finite-dimensional Gaussian
measure naturally embedded in an infinite-dimensional context. Whether this heuristic
connection can be made rigorous is an open problem.

The proposed name for the unified theory --- \textit{Timewave Zero} --- reflects the
three central features: it is a wave ($X(T)$, the band-limited signal); it involves
time (the cumulative phase $\psi$ that reparametrizes the height parameter on the
critical line into the natural clock of arithmetic); and it concerns zeros (the
nontrivial zeros of $\zeta$, which are the transversal nodes of the standing wave).

\bigskip
\hrule
\bigskip

\section{Where the Argument Stands}

The proof as assembled has the following components, each of which is explicit and in
principle verifiable.
\begin{enumerate}
    \item \textbf{$\psi$ is strictly increasing} --- guaranteed by the analyticity of
          $\zeta$ and Ingham's mean-value theorem, which prevents $\psi'$ from
          vanishing on a set of positive measure.
    \item \textbf{$X(T) \in \mathrm{PW}_{[-2,0]}$} --- the core band-limitation claim,
          argued via the Riemann--Siegel expansion and Paley--Wiener--Bernstein. This
          is the step that would require the most careful scrutiny in a formal review,
          as it involves controlling remainder terms in the Riemann--Siegel formula
          under the time change.
    \item \textbf{Hardy space membership $\to$ zero-free half-planes} --- standard
          functional analysis.
    \item \textbf{Schwarz reflection $\to$ all zeros real} --- standard.
    \item \textbf{Bulinskaya $\to$ all zeros simple} --- standard, given
          band-limitation.
\end{enumerate}

The completeness of the argument at step~2 --- specifically the uniform control of the
Riemann--Siegel remainder $R(t)$ under the non-uniform time change $\psi$ --- is the
key technical question for any independent verification. The remaining steps follow from
it by established theorems.

\end{document}
